\documentclass{article}
\usepackage{hyperref}
\usepackage{geometry}
\geometry{margin=1in}

\title{AutoCTM Software Design Specification}
\author{}
\date{}

\begin{document}
\maketitle
\tableofcontents
\newpage

\section{Introduction}

\subsection{Purpose}
This document describes the software design of AutoCTM, a Certificate
Transparency monitor written in Go. It is intended to guide implementation
and serve as a reference for the development team.

\subsection{Scope}
AutoCTM allows operators to monitor one or more RFC 9162-compliant CT logs
for certificate issuance against a configurable domain watchlist. The system
consists of three decoupled components: a CLI, a persistent broker daemon,
and one or more independent monitor instances. The CLI may be closed at any
time without affecting running instances. On relaunch the CLI reconnects to
the broker and resumes control.

\subsection{Definitions}
\begin{itemize}
    \item \textbf{CLI} --- The operator-facing command-line interface. A
    short-lived process that connects to the broker, issues commands, and
    may be closed freely.
    \item \textbf{Broker} --- A persistent background daemon that owns the
    control socket, manages the instance registry, and routes commands from
    the CLI to the appropriate instance.
    \item \textbf{Instance} --- A long-running background process responsible
    for monitoring one or more CT logs, checking certificates, and raising
    alerts.
    \item \textbf{Active Instance} --- The instance currently selected in the
    CLI via \texttt{/set-instance}. Commands issued in instance context are
    routed to this instance.
    \item \textbf{STH} --- Signed Tree Head, as defined in RFC 9162.
    \item \textbf{SCT} --- Signed Certificate Timestamp, as defined in RFC 9162.
    \item \textbf{MMD} --- Maximum Merge Delay, as defined in RFC 9162.
    \item \textbf{Watchlist} --- The set of domains configured by the operator
    for alert monitoring.
\end{itemize}

\subsection{References}
\begin{itemize}
    \item RFC 9162 --- Certificate Transparency Version 2.0
    \item AutoCTM Software Requirements Specification
\end{itemize}

\section{System Overview}

AutoCTM is composed of three decoupled processes:

\begin{itemize}
    \item \textbf{CLI} (\texttt{autoctm}) --- The operator entry point. On
    launch it checks for a running broker. If none is found it spawns one
    automatically before presenting the command prompt. The CLI may be closed
    at any time; running instances are unaffected.
    \item \textbf{Broker} (\texttt{autoctm-broker}) --- A persistent daemon
    that starts automatically on first CLI launch. It owns a Unix socket at a
    known path, maintains the instance registry, and routes commands between
    the CLI and instances. It remains alive independently of both the CLI and
    any instance.
    \item \textbf{Monitor Instance} (\texttt{autoctm-instance}) --- A
    long-running background process spawned by the broker on \texttt{/start}.
    Each instance independently polls its configured CT logs, processes
    entries, performs watchlist and CRL checks, and raises alerts. Multiple
    instances may run concurrently.
\end{itemize}


\section{Architecture}

\subsection{Component Diagram}
% Insert component diagram here

\subsection{Process Lifecycle}

\subsubsection{Broker Startup}
\begin{enumerate}
    \item CLI launches and checks for a broker socket at the well-known path
    (e.g. \texttt{/var/run/autoctm/broker.sock}).
    \item If no socket is found, the CLI forks and execs \texttt{autoctm-broker}
    as a detached background process, then waits for the socket to appear.
    \item The broker creates the socket, loads the instance registry from the
    database, and begins listening for CLI and instance connections.
\end{enumerate}

\subsubsection{CLI Lifecycle}
\begin{enumerate}
    \item On launch, the CLI connects to the broker socket.
    \item The CLI enters a REPL, accepting operator commands and forwarding
    them to the broker.
    \item The CLI maintains an \textit{active instance} context, initially
    unset (global context).
    \item The CLI may be closed at any time with no effect on the broker or
    any running instance. On relaunch it reconnects to the existing broker.
\end{enumerate}

\subsubsection{Instance Lifecycle}
\begin{enumerate}
    \item The broker spawns a new instance process on receipt of \texttt{/start}.
    \item The instance registers itself with the broker, providing its PID and
    a unique instance ID (UUID).
    \item The instance loads its configuration from the database and begins
    its polling loop.
    \item On \texttt{/shutdown} (targeted or global), the broker signals the
    instance to drain its current work and exit cleanly.
    \item The instance deregisters from the broker on exit.
\end{enumerate}

\subsection{Inter-Process Communication}
All IPC uses a Unix domain socket owned by the broker at a well-known path.
Messages are encoded as newline-delimited JSON. There are two connection types:

\begin{itemize}
    \item \textbf{CLI $\rightarrow$ Broker} --- The CLI holds a persistent
    connection to the broker for the duration of its session.
    \item \textbf{Broker $\rightarrow$ Instance} --- The broker holds a
    persistent connection to each registered instance, used to forward
    commands and receive status updates.
\end{itemize}

% ─────────────────────────────────────────────
\section{Component Design}

\subsection{CLI}

\subsubsection{Command Contexts}
The CLI operates in one of two contexts at any time:

\paragraph{Global Context}
No active instance is selected. Commands affect shared state or the broker:
\begin{itemize}
    \item \texttt{/start} --- Instruct the broker to spawn a new monitor
    instance. Prints the assigned instance ID.
    \item \texttt{/get-instances} --- List all registered instances, their
    IDs, and current status.
    \item \texttt{/set-instance <instanceID>} --- Set the active instance.
    The CLI prompt updates to reflect the selected instance. Subsequent
    commands are routed to that instance.
    \item \texttt{/add <domain>} --- Add a domain to the shared watchlist.
    \item \texttt{/remove <domain>} --- Remove a domain from the shared
    watchlist.
    \item \texttt{/shutdown} --- Gracefully shut down all instances and the
    broker. The CLI then exits.
\end{itemize}

\paragraph{Instance Context}
An active instance is selected via \texttt{/set-instance}. Commands are
routed to that instance:
\begin{itemize}
    \item \texttt{/add-log <url>} --- Add a CT log to the instance's
    configuration.
    \item \texttt{/remove-log <url>} --- Remove a CT log from the instance's
    configuration.
    \item \texttt{/pause} --- Pause the instance's polling loop.
    \item \texttt{/resume} --- Resume a paused instance.
    \item \texttt{/status} --- Print the instance's current status, including
    configured logs, last known STH per log, and alert counts.
    \item \texttt{/shutdown} --- Shut down the active instance only.
    \item \texttt{/clear-instance} --- Return to global context.
\end{itemize}

\subsubsection{Auto-Spawn Behaviour}
If the CLI cannot connect to the broker socket on launch it shall
automatically spawn the broker as a detached background process before
presenting the prompt. This requires no manual daemon management from the
operator.

\subsection{Broker}

\subsubsection{Responsibilities}
\begin{itemize}
    \item Own and maintain the Unix control socket.
    \item Maintain the instance registry (running instances, their IDs, PIDs,
    and socket connections).
    \item Route CLI commands to the correct instance.
    \item Spawn new instance processes on \texttt{/start}.
    \item Persist instance registry state to the database so it survives a
    broker restart.
\end{itemize}

\subsubsection{Broker Restart Behaviour}
If the broker is restarted (e.g. after a system reboot), it shall reload the
instance registry from the database and attempt to reconnect to any instances
whose PID is still active. Instances that are no longer running shall be
marked as stopped.

\subsection{Monitor Instance}

\subsubsection{Startup Sequence}
\begin{enumerate}
    \item Receive configuration from the broker (instance ID, initial log
    list if any).
    \item Connect to PostgreSQL and verify schema.
    \item Register with the broker via the control socket.
    \item Spawn one goroutine per configured log.
    \item Enter polling loop.
\end{enumerate}

\subsubsection{Polling Loop}
Each log goroutine shall independently perform the following steps:
\begin{enumerate}
    \item Fetch the current STH via \texttt{GET /ct/v2/get-sth}.
    \item Verify the STH signature using the log's public key.
    \item If the STH is unchanged from the stored STH, wait for the
    configured poll interval and retry.
    \item Fetch new entries since the last known STH via
    \texttt{GET /ct/v2/get-entries}.
    \item For each entry, perform watchlist and CRL checks.
    \item Verify the Merkle consistency proof between the previous and new
    STH via \texttt{GET /ct/v2/get-sth-consistency}.
    \item Persist new entries and update the stored STH in the database.
    \item Dispatch any pending alerts.
    \item Wait for the configured poll interval and return to step 1.
\end{enumerate}

\subsubsection{Watchlist Checking}
Each certificate's Subject CN and SAN fields are matched against all domains
in the watchlist. Matching shall support:
\begin{itemize}
    \item Exact domain match (e.g. \texttt{example.com}).
    \item Wildcard SAN match (e.g. \texttt{*.example.com}).
\end{itemize}
A match shall produce an alert record of type \texttt{watchlist\_match}.

\subsubsection{CRL Checking}
The instance shall maintain a local in-memory cache of CRLs keyed by issuer
DN. CRLs shall be refreshed on a configurable interval. Each observed
certificate serial number shall be checked against the cached CRL for its
issuer. A hit shall produce an alert record of type \texttt{revoked}.


\section{Data Design}
\label{sec:data}

\subsection{Database}
AutoCTM uses a local PostgreSQL instance as its storage backend. The database
is shared across the broker and all monitor instances.

\subsection{Schema}

\subsubsection{certificates}
Stores all observed certificate entries.
\begin{itemize}
    \item \texttt{id} --- Primary key (UUID).
    \item \texttt{leaf\_hash} --- SHA-256 leaf hash (unique).
    \item \texttt{log\_id} --- OID of the source log.
    \item \texttt{observed\_at} --- Timestamp of observation.
    \item \texttt{not\_before}, \texttt{not\_after} --- Certificate validity
    period.
    \item \texttt{subject} --- Certificate subject CN.
    \item \texttt{sans} --- Array of Subject Alternative Names.
    \item \texttt{issuer} --- Issuer DN.
    \item \texttt{serial} --- Certificate serial number.
    \item \texttt{revoked} --- Boolean, updated by CRL checks.
\end{itemize}

\subsubsection{watchlist}
\begin{itemize}
    \item \texttt{id} --- Primary key (UUID).
    \item \texttt{domain} --- Domain pattern (exact or wildcard).
    \item \texttt{added\_at} --- Timestamp of addition.
\end{itemize}

\subsubsection{instances}
\begin{itemize}
    \item \texttt{id} --- Instance ID (UUID).
    \item \texttt{pid} --- OS process ID.
    \item \texttt{status} --- Current status (\texttt{running},
    \texttt{paused}, \texttt{stopped}).
    \item \texttt{started\_at} --- Timestamp.
    \item \texttt{config} --- JSONB blob of instance configuration (log URLs, public keys, poll interval).
\end{itemize}

\subsubsection{sth\_state}
Tracks the last known STH per log per instance.
\begin{itemize}
    \item \texttt{instance\_id} --- Foreign key to \texttt{instances}.
    \item \texttt{log\_id} --- OID of the log.
    \item \texttt{tree\_size} --- Tree size of the last known STH.
    \item \texttt{root\_hash} --- Root hash of the last known STH.
    \item \texttt{timestamp} --- STH timestamp.
\end{itemize}

\subsubsection{alerts}
\begin{itemize}
    \item \texttt{id} --- Primary key (UUID).
    \item \texttt{certificate\_id} --- Foreign key to \texttt{certificates}.
    \item \texttt{instance\_id} --- Foreign key to \texttt{instances}.
    \item \texttt{reason} --- Alert reason (\texttt{watchlist\_match},
    \texttt{revoked}, \texttt{log\_misbehaviour}).
    \item \texttt{detail} --- JSONB blob of alert-specific detail.
    \item \texttt{raised\_at} --- Timestamp.
    \item \texttt{dispatched} --- Boolean, whether the notification was sent.
\end{itemize}

\section{Error Handling}

\subsection{Log Unavailability}
If a CT log is unreachable the instance shall log the failure and back off using exponential retry. If the log remains unavailable beyond the MMD the instance shall raise a \texttt{log\_misbehaviour} alert.

\subsection{Consistency Proof Failure}
A consistency proof verification failure shall be treated as log misbehaviour. The instance shall record the failure, raise an alert, and pause polling of the affected log pending operator review via
\texttt{/resume}.

\subsection{Database Connectivity}
If the database becomes unreachable the instance shall pause processing and retry on a backoff schedule. On recovery the instance shall resume from the last persisted STH state, re-fetching any entries missed during the outage.

\subsection{Broker Unavailability}
If the CLI cannot reach the broker it shall attempt to auto-spawn a newbroker instance. If the broker cannot be spawned the CLI shall print a diagnostic error and exit.

\subsection{Instance Crash}
If an instance exits unexpectedly the broker shall detect the lost connection, mark the instance as \texttt{stopped} in the database, and notify the operator on the next CLI connection.


\end{document}